%%============================ Compiler Directives =======================%%
%%                                                                        %%
% !TeX program = lualatex
% !TeX encoding = utf8
% !TeX spellcheck = russian-aot
%%                                                                        %%
%%============================== Клас документа ==========================%%
%%                                                                        %%
\documentclass[]{article}
\usepackage[fontsize = 12pt]{fontsize}
\usepackage{ifluatex}
%%                                                                        %%
%%========================== Мови, шрифти та кодування ===================%%
%%
\ifluatex                                                                        %%
	\usepackage{fontspec}
	\setsansfont{CMU Sans Serif}%{Arial}
	\setmainfont{CMU Serif}%{Times New Roman}
	\setmonofont{CMU Typewriter Text}%{Consolas}
	\defaultfontfeatures{Ligatures={TeX}}
	\usepackage[math-style=TeX]{unicode-math}
\else
	\usepackage[utf8]{inputenc}
	\usepackage[T2A,T1]{fontenc}
	\usepackage{amsmath}
	%\usepackage{pscyr}
	\usepackage{cmap}
\fi
\usepackage[english, russian]{babel}
%%                                                                        %%
%%============================= Геометрія сторінки =======================%%
%%                                                                        %%
\usepackage[%
	a4paper,%
	footskip=1cm,%
	headsep=0.3cm,%
	top=2cm, %поле сверху
	bottom=2cm, %поле снизу
	left=2cm, %поле ліворуч
	right=2cm, %поле праворуч
    ]{geometry}
%%                                                                        %%
%%============================== Інтерліньяж  ============================%%
%%                                                                        %%
\renewcommand{\baselinestretch}{1}
%-------------------------  Подавление висячих строк  --------------------%%
\clubpenalty =10000
\widowpenalty=10000
%---------------------------------Інтервали-------------------------------%%
\setlength{\parskip}{0.5ex}%
\setlength{\parindent}{2.5em}%
%%                                                                        %%
%%=========================== Математичні пакети і графіка ===============%%
%%                                                                        %%
\usepackage{tikz, pgfplots}
\usetikzlibrary{calc}

\usepackage{tcolorbox}
%%                                                                        %%
%%========================== Гіперпосилення (href) =======================%%
%%                                                                        %%
\usepackage[%colorlinks=true,
	%urlcolor = blue, %Colour for external hyperlinks
	%linkcolor  = malina, %Colour of internal links
	%citecolor  = green, %Colour of citations
	bookmarks = true,
	bookmarksnumbered=true,
	unicode,
	linktoc = all,
	hypertexnames=false,
	pdftoolbar=false,
	pdfpagelayout=TwoPageRight,
	pdfauthor={Ponomarenko S.M. aka sergiokapone},
	pdfdisplaydoctitle=true,
	pdfencoding=auto
	]%
	{hyperref}
		\makeatletter
	\AtBeginDocument{
	\hypersetup{
		pdfinfo={
		Title={\@title},
		}
	}
	}
	\makeatother
\title{Космология}
%\author{С.Н. Пономаренко}
\date{}

\begin{document}
\maketitle

\section{Кинематика Вселенной}

Однородная и изотропная Вселенная может быть описана нестационарной (т.е. зависящей от времени) метрикой специального вида (т.н. метрика Фридмана-Робертсона-Уокера)

Для понимания того, как может выглядеть наша трёхмерная Вселенная (если она не плоская, конечно), там нужно <<посмотреть>> на неё с точки зрения четвёртого измерения.

Допустим это четырехмерная гиперсфера радиуса $a$ (на поверхности которой мы живём и не ведаем про $x_4$):
\begin{equation*}\label{}
    x_1^2 + x_2^2 + x_3^3 + x_4^2 = a^2
\end{equation*}

Элемент длины на поверхности такой сферы в декартовых координатах:
\begin{equation*}\label{}
    dl^2 = dx_1^2 + dx_2^2 + dx_3^3 + dx_4^2.
\end{equation*}

Исключим из этого выражения <<фиктивную>> координату $x_4$:
\begin{equation}\label{}
    dl^2 = dx_1^2 + dx_2^2 + dx_3^3 + \frac{(x_1dx_1 + x_2dx_2 + x_3dx_3)^2}{a^2 - x_1^2 - x_2^2 - x_3^3}.
\end{equation}

Теперь введём сферические координаты $r, \theta, \phi$, откуда получим:
\begin{equation*}\label{}
	dl^2=  \frac{dr^2}{1-\frac{r^2}{a^2}}+r^2\left[ d\theta^2+\sin^2\theta d\phi^2\right] ,
\end{equation*}

Аналогично можно поступить для плоский и гиперболической геометрии и получить квадрат элемента длины в общем виде:
\begin{equation*}\label{}
	dl^2=  \frac{dr^2}{1-k\frac{r^2}{a^2}}+r^2\left[ d\theta^2+\sin^2\theta d\phi^2\right] ,
\end{equation*}
где постоянная $k$ определяет одну из трёх возможных глобальных топологий пространства:
\begin{enumerate}
\item плоское, $k=0$,
\item постоянной положительной кривизны, $k=+1$,
\item постоянной отрицательной кривизны, $k=-1$.
\end{enumerate}

Пространственно-временная метрика будет иметь вид:
\begin{equation}\label{eq:FWL}
	ds^2=c^2dt^2-\frac{dr^2}{1-\frac{r^2}{a^2}}+r^2\left[ d\theta^2+\sin^2\theta d\phi^2\right],
\end{equation}

В конформной пространственной метрике:
\begin{equation}\label{eq:FWL_conf}
    	ds^2=c^2dt^2-a^2(t)\left[  d\chi^2 + f^2(\chi)\left( d\theta^2+\sin^2\theta d\phi^2\right)    \right],
\end{equation}
где
\begin{equation*}\label{}
    f^2(\chi) =
            \begin{cases}
                \chi, \quad \text{для плоской метрики},\\
                \sin\chi, \quad \text{для пространства положительной кривизны метрики},\\
                \sh\chi, \quad \text{для пространства отрицательной кривизны метрики}.
            \end{cases}
\end{equation*}
$a(t)$ -- Масштабный фактор, единственная зависящая от времени величина, к тому же имеет размерность расстояния. Конформная метрика интересна тем, что она похожа на статичную метрику.

\subsection{Закон Хаббла и красное смещение}

\begin{center}
\begin{tikzpicture}[
    declare function = {
            xm(\n) = sqrt(2)*sqrt(2*pi*\n+pi+2) - 2;
            f(\x) = sin(deg(\x*(1+0.25*\x)));
        },
    ]
    \draw[color=blue, samples=500] plot[domain=3.45:15.11] (\x,{f(\x)});
    \draw ({xm(2)}, 1) -- ++(0, 0.5) ({xm(4)}, 1) -- ++(0, 0.5);
    \draw[latex-latex] ({xm(2)}, 1.25) -- node[above] {$\lambda_0$} ({xm(4)}, 1.25);
    \draw ({xm(20)}, 1) -- ++(0, 0.5) ({xm(22)}, 1) -- ++(0, 0.5);
    \draw[latex-latex] ({xm(20)}, 1.25) -- node[above] {$\lambda$} ({xm(22)}, 1.25);
\end{tikzpicture}
\end{center}
С течением времени пространство расширяется, и длина волны фотона увеличивается вместе с расширением пространства. Если в момент времени $t$ испускается фотон с длиной волны $\lambda$ то уже в современный момент времени $t_0 $ он уже будет приниматься с длиной волны $\lambda_0$:
\begin{equation}\label{}
	\lambda_0 = \frac{a_0}{a(t)}\lambda.
\end{equation}
Величину, которую обозначают через $z$ ---  называют красным смещением:
\begin{equation}\label{key}
	\frac{a_0}{a(t)} = 1 + z(t).
\end{equation}
Как раз наблюдая спектры удалённых галактик определяют именно $z$.

При $z \ll 1$, можно разложить $a(t)$ в ряд Тейлора
\begin{equation*}\label{}
	a(t) = a_0 - \dot{a} (t_0 - t)
\end{equation*}
и тогда
\begin{equation*}\label{}
	\frac{a_0}{a(t)} = 1 + z(t) = \frac{a_0}{a_0 - \dot{a}_0 (t_0 - t)} = \frac{1}{1 - \frac{\dot{a}_0}{a_0}(t_0 - t)} = 1 +  \frac{\dot{a}}{a}(t_0 - t).
\end{equation*}
Введём параметр Хаббла как:
\begin{equation}\label{eq:H}
	H_0 = \frac{\dot{a}_0}{a_0},
\end{equation}
и тогда При $z \ll 1$ получаем
\begin{equation}\label{eq:low_z}
	z(t) = H_0(t_0 - t).
\end{equation}
Поскольку в наших единицах $c = 1$, то $c(t_0 - t) = r$ есть расстояние между источником и приёмником света, т. е.
\begin{equation}\label{eq:Habbl_Law}
	z(t) = H_0r.
\end{equation}
Эта формула и называется законом Хаббла. По результатам Plank $H_0 = 67.8 \pm 0.9\, \frac{\text{км/с}}{\text{МПк}}$.
Кстати $z$ есть скорость источника света, выраженная в единицах скорости света.

\begin{tcolorbox}
	Какова скорость галактики, красное смещение которой $z = 0.2$? Какое расстояние до этой галактики?
\end{tcolorbox}

\begin{tcolorbox}
	Найти первую поправку к формуле \eqref{eq:Habbl_Law}. Выразить ее через параметр ускорения $q_0 = \frac{1}{H_0^2} \left. \frac{\ddot{a}}{a} \right|_{t = t_0} $
\end{tcolorbox}

\begin{tcolorbox}[title = Замена переменных $a \to z$. Бокс \ref{box:az}]\label{box:az}
Часто нам будет необходима замена переменных $a \to z$.
Поскольку $a_0/a = 1 + z$ то $da/a = -dz/(1+z)$.
\end{tcolorbox}
\subsection{Расстояния во Вселенной}

\subsubsection{Фотометрическое расстояние}
Фотометрическое расстояние связано с измерением яркости удалённых объектов и определением расстояния по Этой яркости.

Введём понятия:
\begin{enumerate}
\item $L$~--- абсолютная светимость --- энергия, которую излучает объект за единицу времени [эрг/c].
\item $J$~--- видимая яркость объекта --- плотность потока энергии принимаемая прибором [эрг/с/м$^2$].
\end{enumerate}
В плоской нерасширяющейся Вселенной простое соотношение между этими величинами:
\begin{equation*}\label{}
    J = \frac{L}{4\pi r^2},
\end{equation*}
где $r$ --- расстояние до источника.

В расширяющейся Вселенной, расстояние $r = r_\text{ph}$~--- называется фотометрическим. Выясним, как будет зависеть $r_\text{ph}$ от красного смещения и от геометрии Вселенной.
\begin{equation*}\label{eq:J}
    J = \frac{L}{4\pi r_\text{ph}^2}.
\end{equation*}

На данный момент площадь сферы, на которую разошлось излучение удалённого источника с координатой $\chi$ равна:
\begin{equation*}\label{}
    S(\chi) = 4\pi a_0^2 f^2(\chi).
\end{equation*}
Это расстояние $a_0 f(\chi)$ пока не фотометрическое расстояние так как нужно учесть то, что излучение <<покраснело>> с момента излучения, а значит его энергия, которая до нас дошла уменьшилась в $1 + z$ раз, т. е. $L \to \frac{L}{1 + z}$.

Кроме того, поток энергии, воспринимаемый прибором тоже уменьшается в $1 + z$ раз, т. е. $J \to \frac{J}{1 + z}$.

Поэтому фотометрическое расстояние:
\begin{equation}\label{eq:ph_dist}
    r_\text{ph} = (1+z) a_0 f(\chi(z)),
\end{equation}
где имеется ввиду, что $z$~--- красное смещение источника, испустившего сигнал.

Для определения $r_\text{ph} (z)$ нам нужно ещё выразить $\chi(z)$. Из метрики \eqref{eq:FWL_conf} положив $ds^2 = 0$ для света находим:
\begin{equation*}\label{}
    \chi = \int\limits_t^{t_0}\frac{dt}{a(t)} = \int\limits_t^{t_0} \frac{da}{H(a)a^2} \stackrel{\ref{box:az}}{=} \int\limits_0^{z} \frac{dz}{a_0 H(z)},
\end{equation*}
выделим последнюю формулу отдельно:
\begin{equation}\label{eq:chi_z}
    \chi(z) = \int\limits_0^{z} \frac{dz}{a_0 H(z)}.
\end{equation}

\subsection{Расстояние углового размера}

Расстояние углового размер связано с понятием угла, под которым видны удалённые протяжённые объекты поперечное физическое расстояние  между которыми равно $l = a f(\chi) \delta\theta$:
\begin{equation*}\label{}
    \delta\theta = \frac{l}{r_A},
\end{equation*}
где $r_A$ --- расстояние углового размера.


\begin{center}
\begin{tikzpicture}
    \def\p{2}
    \draw (0,0) circle (\p) circle (\p);
    \draw[] (0,0) -- node[right] {$r_A$} (75:\p) (0,0) -- (105:\p);
    \draw[red, ultra thick] (0,0) ++(105:\p) arc (105:75:\p) node[above, pos=0.5, color=black] {$l$};
    \def\p{4}
    \draw (0,0) circle (\p) circle (\p);
    \draw[] (0,0) --  (75:\p) (0,0) -- (105:\p);
    \draw[red, ultra thick] (0,0) ++(105:\p) arc (105:75:\p) node[above, pos=0.5, color = black] {$a_0f(\chi)\delta\theta$};
    \def\p{0.75}
    \draw[] (0,0) ++(105:\p) arc (105:75:\p) node[above, pos=0.5] {$\theta$};
\end{tikzpicture}
\end{center}

Расстояние углового размера равно
\[
    r_A = a f(\chi) = \frac{a_0 f(\chi)}{1 + z} = \frac{r(z)}{1 + z},
\]
где $r(z)$ --- расстояние до объектов с красным смещением $z$ в современную эпоху, тогда как $r_A$ --- расстояние до объектов в момент испускания света (заметим, что $f(\chi)$ --- и в момент испускания и в современную эпоху одинаковы).


\section{Динамика Вселенной}

Динамика Вселенной должна ответить на вопрос, как зависит масштабный фактор $a(t)$ от времени и чем определяется эта зависимость.

\subsection{Уравнения Эйнштейна и Фридмана}

Логично использовать для этого уравнения гравитационного поля Эйнштейна используя метрику \eqref{eq:FWL}.
Уравнения Эйнштейна принимают вид:
\begin{equation}\label{eq:Freiedman}
	\left( \frac{\dot{a}}{a}\right) = \frac{8\pi}{3}G \rho - \frac{k}{a^2}.
\end{equation}

Это уравнение называют уравнением Фридмана; оно связывает темп расширения Вселенной (а именно, параметр Хаббла $H = \dot{a}/a$) с плотностью энергии материи $\rho$ и пространственной кривизной.

Уравнение Фридмана \eqref{eq:Freiedman} содержат две неизвестные функции времени $a(t)$ и $\rho(t)$.
Для получения дополнительного уравнения удобно рассмотреть условие ковариантного сохранения тензора энергии-импульса вещества $\nabla_{\mu} T^{\mu\nu} = 0$, которое даёт:
\begin{equation}\label{eq:rho}
	\dot{\rho}(t) + 3H(\rho + p) = 0.
\end{equation}

Для замыкания системы уравнений, определяющих динамику эволюции однородной изотропной Вселенной, необходимо задать ещё уравнение состояния материи:
\begin{equation}\label{eq:state}
	p = p(\rho).
\end{equation}

Уравнения \eqref{eq:Freiedman}, \eqref{eq:rho} и \eqref{eq:state} полностью определяют динамику расширения Вселенной.

\subsection{Обезразмеренное уравнение Фридмана и критическая плотность}

Проанализируем уравнение \eqref{eq:Freiedman}. Для этого поделим его на $H^2$  и запишем обезразмеренном виде:
\begin{equation}\label{eq:1}
	\Omega_{\rho} = 1- \Omega_{K},
\end{equation}
где $\Omega_{\rho} = \dfrac{8\pi G \rho}{3H^2}$ --- отражает вклад плотности Вселенной, $\Omega_{k} = \dfrac{k}{a^2H^2}$ --- вклад кривизны Вселенной.

Во всей этой истории, мы еще не знаем топологию пространства, т.е. мы не знаем $k$. Поскольку
\begin{equation*}\label{}
	\rho = \frac{3H^2}{8\pi G} \Omega_{\rho},
\end{equation*}
то можно ввести понятие критической плотности, положив $\Omega_{\rho} = 1$ (тогда уравнение \eqref{eq:1} выполнится когда $\Omega_k = 0$):
\begin{equation}\label{eq:cr}
	\rho_c = \frac{3H^2}{8\pi G},
\end{equation}
т. е. критическая плотность, это такая плотность, при которой реализуется плоская Вселенная.

И теперь, согласно \eqref{eq:1}  если
\begin{itemize}
	\item $\rho < \rho_c$ --- у нас будет Вселенная с отрицательной кривизной --- открытая Вселенная;
	\item $\rho = \rho_c$ --- у нас будет Вселенная с нулевой кривизной --- плоская Вселенная;
	\item $\rho > \rho_c$ --- у нас будет Вселенная с положительной кривизной --- замкнутая Вселенная.
\end{itemize}


\subsection{Что и как можно узнать из наблюдений?}

\begin{tcolorbox}
	Как узнать чему равна критическая плотность Вселенной?
\end{tcolorbox}
Ответ: узнайте параметр Хаббла и подставьте в \eqref{eq:cr}.

Наблюдательные данные по определению параметра Хаббла дают оценку:
\[
	\rho_c = 9.31\cdot 10^{-27}
	\, \text{кг/м${}^3$}
\]
Критическая плотность соответствует $5.5$ атома водорода на кубический метр.


\begin{tcolorbox}
	Вопрос: как узнать чему равна плотность Вселенной?
\end{tcolorbox}

Узнайте $\Omega_{\rho}$ и умножьте на критическую плотность $\rho = \rho_c \Omega_{\rho}$.

\begin{tcolorbox}
	Как узнать чему равна $\Omega_{\rho}$?
\end{tcolorbox}

По результатам WMAP
\begin{equation}\label{}
	0.9915 < \Omega_{\rho}^{(0)} <  1.0085,
\end{equation}
т.е. для $\Omega_{k}^{(0)}$ ($(0)$ --- означает <<в современную эпоху>>) из \eqref{eq:1} получается:
\begin{equation}\label{}
	-0.0175 < \Omega_{k}^{(0)} <  0.0085.
\end{equation}
Т. е. плотность во Вселенной с высокой точностью равна критической, а значит, мы живём в плоской Вселенной.

\bigskip

Вклады в $\Omega_{\rho}^{(0)} = \Omega_{M}^{(0)} + \Omega_\text{rad}^{(0)} + \Omega_{\Lambda}^{(0)} $:
\begin{enumerate}
	\item Вклад материи (включая темную и барионную) $\Omega_{M}^{(0)} = 0.31 \pm 0.01$ (барионная материя, которая входит в эту цифру $\Omega_B = 0.048 \pm 0.002$).
	\item Вклад излучения $\Omega_\text{rad}^{(0)}$: реликтовые фотоны $\Omega_{\gamma}^{(0)} \approx 5\cdot 10^{-5}$ и нейтрино $\Omega_{\nu}^{(0)} \approx 3.5\cdot 10^{-5}$.
	\item Вклад тёмной энергии $\Omega_{\Lambda}^{(0)} = 0.69 \pm 0.01$.
\end{enumerate}

\begin{tcolorbox}
	Как по результатам WMAP (Plank) определили $\Omega_{\rho}$ и вклады?
\end{tcolorbox}


\subsection{Параметризированное уравнение Фридмана}

Уравнение Фридмана \eqref{eq:Freiedman} можно записать таким образом, чтобы оно включало себя современные значения $H_{0}^2$, $\Omega$ и $a_0$:
\begin{equation}\label{eq:Friedman_Parametric}
	H^2 = H_0^2\left[ \Omega_{M}^{(0)} \left( \frac{a_0}{a}\right)^3 + \Omega_{\Lambda}^{(0)} + \Omega_\text{rad}^{(0)}\left( \frac{a_0}{a}\right)^4 + \Omega_{k}^{(0)} \left( \frac{a_0}{a}\right)^2  \right],
\end{equation}
Кстати согласно $\Lambda\mathrm{CDM}$ $\Omega_{\Lambda}^{(0)} = \Omega_{\Lambda} = \mathrm{const}$.

\begin{tcolorbox}
	При каком $z$ плотности материи и тёмной энергии совпадали?
\end{tcolorbox}
\begin{equation*}\label{}
	\Omega_{M}^{(0)} \left( \frac{a_0}{a}\right)^3 = \Omega_{\Lambda}.
\end{equation*}
Отсюда $z = 0.3$.
\begin{tcolorbox}
	При каком $z$ замедленное расширение перешло в ускоренное?
\end{tcolorbox}
Нужно найти $\ddot{a}$ и приравнять 0.
Из \eqref{eq:Friedman_Parametric} (эпоха доминирования радиации закончилась, а кривизной пренебрегаем):
\begin{equation*}
	\left( \frac{\dot{a}}{a}\right)^2 = H_0^2\left[ \Omega_{M}^{(0)} \left( \frac{a_0}{a}\right)^3 + \Omega_{\Lambda}^{(0)} \right],
\end{equation*}
откуда умножив слева и справа на $a$ и продифференцировав по $t$:
\begin{equation*}
	2\ddot{a}\dot{a} = H_0^2\left( -\Omega_{M}^{(0)}  \frac{a_0^3}{a}  + 2\Omega_{\Lambda} a \right) \dot{a}.
\end{equation*}
\begin{equation*}
	1+ z = \left( \frac{2\Omega_{\Lambda}}{\Omega_{M}^{(0)}}\right)^{\kern-5pt\frac13}.
\end{equation*}
Откуда $z = 0.67$.

\begin{tcolorbox}
	При каком $z$ плотности радиации и материи совпадали?
\end{tcolorbox}
Из \eqref{eq:Friedman_Parametric} ($\Omega_{\Lambda}$ пренебрегаем):
\begin{equation}\label{}
	1 + z = \frac{\Omega_{M}}{\Omega_\text{rad}},
\end{equation}
откуда $z = 3\,600$.
Это важный момент, поскольку до этого произошло изменение эволюции неоднородностей плотности материи. С этого момента неоднородности стали такими, что смогли образоваться галактики, скопления \ldots

\subsubsection{Оценка времён}

Решаем уравнение~\eqref{eq:Friedman_Parametric} относительно $z$, предварительно заменив $a_0/a = z$, $da/a = -dz/(1+z)$ и разделив переменные:
\begin{equation}\label{eq:tz}
 t(z) = \frac1{H_0}\int\limits_z^{\infty} \frac{dz}{(1+ z)\sqrt{\Omega_{M}^{(0)} (1 + z)^3 + \Omega_\text{rad}^{(0)}(1 + z)^4  + \Omega_{\Lambda}}}
\end{equation}
Этот интеграл будем брать с помощью \texttt{Wolfram Alpha}:
{\footnotesize
\begin{verbatim}
    3.2e-8*(1/2.2e-18)*Integrate[1/((1 + z) Sqrt[0.31 (1 + z)^3 + 7.5e-5 (1 + z)^4 + 0.69]),
    {z, 0, Infinity}]
\end{verbatim}
}
\noindent где взят параметр Хаббла $H_0 = 2.2\cdot 10^{-18}$~1/с (в обратных секундах), а $3.2\cdot 10^{-8}$~--- коэфициент перевода из секунд в года.

\begin{tcolorbox}
    Оценить возраст Вселенной на момент перехода из стадии доминирования радиации в стадию доминирования материи.
\end{tcolorbox}
Положим \href{https://bit.ly/2Zv1GKE}{$z = 3\,600$}. Получим значение $t \approx 49\,000 $~лет.

\begin{tcolorbox}
    Оценить возраст Вселенной на текущий момент.
\end{tcolorbox}
Для оценки возраста Вселенной на текущий момент, положим \href{https://bit.ly/32wa4uw}{$z = 0$}. Получим значение $t \approx 13.75 $~млрд. лет.

\begin{tcolorbox}
    Получите зависимость возраста вселенной от $z$.
\end{tcolorbox}
%В формуле \eqref{eq:tz} пренебрежём временем до начала эры доминирования материи, тогда интеграл упростится до вида:
%Этот интеграл будем брать с помощью \texttt{Wolfram Alpha}:
%{\footnotesize
%\begin{verbatim}
%    3.2e-8*(1/2.2e-18)*Integrate[1/((1 + z) Sqrt[0.31 (1 + z)^3 + 0.69]),
%    {z, 0, Infinity}]
%\end{verbatim}
%}

\subsection{Инфляция, Большой Взрыв и проблема плоскостности}

Наша Вселенная практически плоская $\Omega_{k}^{(0)} \approx 0$ согласно современным данным, но вклад кривизны убывает как $\Omega_{k} \propto 1/a^2 $, а вклад материи как $\Omega_{M} \propto 1/a^3$ и излучения  $\Omega_\text{rad} \propto 1/a^4$ , т. е. экстраполируя в раннюю эпоху $\Omega_{k}$ должен становится ещё меньшим по сравнению с $\Omega_{M}$ и $\Omega_\text{rad}$.

Вывод: на ранних этапах доминировало в большей степени $\Omega_\text{rad}$ и в меньшей степени $\Omega_{M}$, а $\Omega_{k} \approx 0$ (кривизна вообще не существенна, даже если и $k \neq 0$) по прежнему.

\begin{tcolorbox}
	А откуда тогда взялись $\Omega_\text{rad}$ и $\Omega_{M}$?
\end{tcolorbox}

Скорее всего (гипотеза), было некое поле с постоянной плотностью (инфлатон):
\begin{equation}\label{eq:inflaton}
	\Omega_\text{inf} \approx 1,
\end{equation}
тогда из \eqref{eq:Friedman_Parametric} получим закон изменения $a(t)$:
\begin{equation}\label{eq:at}
	a \sim e^{Ht}
\end{equation}
где $H = 10^{37}$~1/с, т. е. характерное время удвоения размера вселенной $\tau = 1/H = 10^{-37}$~с. Это значит, что из-за такого <<катастрофического>> расширения, Вселенная <<уплощилась>>. А в конце экспоненциального расширения инфлатон распался на вещество, радиацию и оставило огарок в виде тёмной энергии ---  этот распад и называется <<Большой Взрыв>>.



\end{document}


